% --- Mode d'emploi ----------------------------------------------------------
\chapter*{Comment ce livre est fait}
\markboth{Comment ce livre est fait}{Comment ce livre est fait}
\addcontentsline{toc}{chapter}{Comment ce livre est fait}

Tous les chapitres suivent la même route, dans le même ordre. Au bout de deux
ou trois chapitres, vous saurez où regarder sans chercher.

\vspace{0.8em}

\begin{revision}
On rouvre ce que vous savez déjà et dont le chapitre va se servir : une
formule de première, une propriété de seconde. Jamais un mot de la notion
nouvelle.
\end{revision}

\begin{automatismes}
Une dizaine de calculs qui doivent sortir sans papier. C'est le
chronomètre, pas la réflexion, qui compte ici.
\end{automatismes}

\noindent Vient ensuite \textbf{le cours}, où alternent trois formes :

\begin{definition}[Ce qu'on pose]
Le mot nouveau, écrit une fois pour toutes. On y revient chaque fois qu'un
doute apparaît.
\end{definition}

\begin{theoreme}[Ce qu'on démontre]
Le résultat qui fait travailler la définition, suivi de sa preuve.
\end{theoreme}

\begin{demonstration}
La preuve est écrite en petit, en retrait. Elle se lit après le théorème, ou
plus tard, ou jamais si le temps manque — mais elle est là.
\end{demonstration}

\begin{visualisation}[la figure qui montre l'idée]
Une figure, et à côté la phrase qu'un professeur dirait en la traçant au
tableau. Regardez d'abord le dessin, lisez la légende ensuite.
\end{visualisation}

\begin{exemple}[le calcul déroulé en entier]
Un cas traité ligne à ligne, sans étape sautée, avec les justifications dans
la marge.
\end{exemple}

\begin{methode}{reconnaître le bon outil}
Quand plusieurs résultats s'appliquent, ce bloc dit lequel choisir et à quoi
on le reconnaît.
\end{methode}

\begin{erreurclassique}
La faute que l'on corrige chaque année sur les copies, écrite noir sur blanc
pour que vous la voyiez venir.
\end{erreurclassique}

\noindent Puis on passe la main :

\begin{entrainement}[juste après la notion]
Des exercices courts, placés immédiatement après la notion qu'ils font
travailler, classés du plus simple au plus exigeant. Les pastilles
\niveau{1}, \niveau{2}, \niveau{3} donnent le niveau.
\end{entrainement}

\begin{jecherche}[une situation à démêler]
Un problème plus ouvert, posé \emph{après} le cours, où plusieurs chemins
mènent au résultat — souvent ancré dans une question économique ou de
gestion, à l'image de l'épreuve.
\end{jecherche}

\begin{jemeteste}
Un questionnaire de fin de chapitre, à faire seul, pour savoir si le chapitre
est acquis avant de passer au suivant.
\end{jemeteste}

\begin{commeaubac}
Un exercice au format et au barème de l'épreuve : quatre heures, coefficient
cinq, calculatrice scientifique non programmable et table financière
autorisées.
\end{commeaubac}

\begin{notepourlaclasse}
Ce bloc n'apparaît que dans l'édition du professeur : durée conseillée,
difficulté prévisible, variante pour classe à grand effectif.
\end{notepourlaclasse}

\begin{remarque}[une différence avec le livre de série S]
Ce livre ne comporte pas de bloc « je vise le concours » : aucun sujet
d'entrée à l'ENI ou à l'ESPA ne correspond au profil de la série OSE. La
dernière étape de chaque chapitre est donc l'exercice au format du
baccalauréat.
\end{remarque}
